\documentclass[../main.tex]{subfiles}
\graphicspath{{\subfix{../img/}}}
\begin{document}
	\subsection{Corrente Elettrica Stazionaria}
	N.B: L'intensità non dipende dal tempo.\\\\
	Negli atomi di un metallo gli elettroni periferici non si legano ai singoli atomi, ma sono liberi di muoversi nel reticolo formato dagli ioni positivi e sono detti elettroni di conduzione.\\
	Nel rame vi è un elettrone libero per atomo e quindi il numero di elettroni di conduzione è di circa $8 \times 10^{22}$ elettroni/cm$^3$.\\\\
	L’ordine di grandezza è lo stesso per tutti i conduttori metallici.
	$$e = -1.6 \times 10^{-19} \text{ C}$$
	$$ m = 9.1 \times 10^{-31} \text{ Kg}$$
	Il moto degli elettroni liberi in un conduttore in equilibrio elettrostatico è disordinato. \\
	In qualsiasi volume piccolo su scala macroscopica ma contenente un numero $N$ di elettroni abbastanza elevato (in un $\mu \, m^3$ si hanno circa $10^{11}$ elettroni), la velocità media è nulla (Che non è altro che una media aritmetica):
	$$\mathbf{v} = \frac{1}{N} \sum_{i=1,N} \mathbf{v}_i = 0$$
	$$v = \sqrt{\frac{3KT}{m}} \approx 10^5 \text{ m/s}$$
	\begin{center}
		\begin{figure}[h]
			\centering
			\includegraphics[scale=0.55]{65.png}
			\caption{}
			\label{fig:im-65}
		\end{figure}
	\end{center}
	Non esiste una direzione di moto preferenziale per gli elettroni.\\
	Occasionalmente un elettrone “urta” uno ione del cristallo subendo una brusca variazione di direzione del moto.
	\begin{center}
		\begin{figure}[h]
			\centering
			\includegraphics[scale=0.55]{66.png}
			\caption{Dove $\tau$ è il tempo medio trascorso tra una collisione e l’altra}
			\label{fig:im-66}
		\end{figure}
	\end{center}
	Quando si applica un campo elettrico $E$ gli elettroni sono sottoposti ad una forza $F = -e \cdot E$ e modificano il loro moto casuale migrando \textbf{lentamente} in verso opposto al campo con una velocità di deriva $v_d$ (tipicamente dell’ordine di $10^{-5}$m/s).
	\begin{center}
		\begin{figure}[h]
			\centering
			\includegraphics[scale=0.55]{67.png}
			\caption{Dove $\tau$ è il tempo medio trascorso tra una collisione e l’altra}
			\label{fig:im-67}
		\end{figure}
	\end{center}
	Il flusso ordinato di cariche elettriche è detto \textbf{corrente elettrica}.
	\begin{center}
		\begin{figure}[h]
			\centering
			\includegraphics[scale=0.55]{68.png}
			\caption{Dove $\tau$ è il tempo medio trascorso tra una collisione e l’altra}
			\label{fig:im-68}
		\end{figure}
	\end{center}
	Quello che va veloce nel circuito è l’informazione di come si propaga il campo elettrico, gli elettroni vanno molto lentamente.\\\\
	Una quantità di carica $dq$ transita nel tempo $dt$ attraverso un’area $S$, ad es. la sezione trasversale di un filo percorso da corrente. \\
	Si definisce \textbf{intensità di corrente} $i$ la carica netta che attraversa tale superficie nell’unità di tempo:
	$$i = \frac{dq}{dt}$$
	Convenzione: il verso della corrente è quello in cui si muovono le cariche positive, anche quando nella realtà i portatori di carica sono negativi.\\\\
	Si ha l’intensità di corrente di $1$A quando, attraverso una data superficie, passa una carica di un Coulomb in un secondo.
	$$[A] = [C/s]$$
	\subsection{Densità di corrente}
	La carica netta che attraversa la superficie in un certo tempo si trova integrando la corrente nel tempo:
	$$q = \int i \, dt$$
	Una grandezza vettoriale legata alla corrente è la densità di corrente, ossia la corrente per unità di area, il cui modulo è ($S$ perpendicolare al filo e velocità dei portatori uniforme su $S$)
	$$j = \frac{i}{S}$$
	La direzione e il verso di $j$ coincidono con quelli del moto delle cariche positive.\\\\
	La corrente che scorre attraverso una generica superficie, suddivisa in elementi di area individuati da $dS$ (ortogonale all’elemento di area e con verso concorde a $j$), si determina integrando $j$ su tutta la superficie:
	$$i = \int_S \mathbf{j} \cdot \mathbf{dS}$$
	L’intensità di corrente attraverso una superficie è il \textbf{flusso della densità di corrente} attraverso quella superficie.
	\begin{center}
		\begin{figure}[h]
			\centering
			\includegraphics[scale=0.55]{69.png}
			\caption{Dove $v_d$ indica la velocità di deriva}
			\label{fig:im-69}
		\end{figure}
	\end{center}
	Se ho $n$ elettroni per unità di volume, percorrono un tratto $L$ nel tempo $t = \frac{L}{v_d}$.\\
	Attraverso la sezione $A$ nel tempo $t$ transitano tutti gli elettroni contenuti nel volume $AL$.
	La carica netta che attraversa la superficie in un tempo $t$ è quindi:
	$$q = -neAL$$
	$$j = \left| \frac{q}{At} \right| = \frac{enAL}{AL/v_d} = e \cdot n \cdot v_d$$
	$$\boxed{j = -e\cdot n \cdot v_d}$$
	\textbf{Ha lo stesso verso di $E$}.\\\\
	In generale se si hanno $n$ portatori di carica $q$ per unità di volume.
	$$\mathbf{j} = n \cdot q \cdot v_d$$
	\textbf{Ha sempre lo stesso verso di $E$}.\\\\
	Se si hanno due tipi di portatori:
	$$\mathbf{j} = n_a q_a v_{da} + n_b q_b v_{db}$$
	\begin{center}
		\begin{figure}[h]
			\centering
			\includegraphics[scale=0.55]{70.png}
			\caption{}
			\label{fig:im-70}
		\end{figure}
	\end{center}
	Per la conservazione della carica elettrica, in condizioni stazionarie l’intensità di corrente è la stessa attraverso ogni sezione del conduttore (la carica non si accumula, la densità di carica è costante).\\
	Se il conduttore è a sezione variabile la densità di corrente, e quindi la velocità di deriva sono maggiori dove la sezione è minore.
	\subsection{Legge dei nodi (1$^a$ legge di Kirchoff)}
	\begin{center}
		\begin{figure}[h]
			\centering
			\includegraphics[scale=0.55]{71.png}
			\caption{$i_1 - i_2 - i_3 + i_4 + i_5 = 0$}
			\label{fig:im-71}
		\end{figure}
	\end{center}
	$$\sum_i i_i = 0$$
	Un nodo è un punto nel quale convergono almeno tre conduttori.\\\\
	La prima legge di Kirchoff afferma che:\\
	In condizioni stazionarie la somma algebrica delle correnti uscenti da un nodo è nulla (per convenzione si prendono come positive le correnti uscenti).
	\subsection{Legge di Ohm}
	Se agli estremi di un tratto di conduttore come un filo metallico viene applicata una differenza di potenziale $\Delta V$, nel conduttore si produrrà una corrente di intensità $i$.\\\\
	L’intensità di corrente che circola in un conduttore è, per molti materiali, direttamente proporzionale alla differenza di potenziale $\Delta V$:
	$$\Delta V = Ri$$
	Di conseguenza è presente un campo elettrico.\\\\
	$R$ è una proprietà del conduttore ed è chiamata \textbf{resistenza} e si misura in Ohm ($\Omega = V/A$).
	\begin{center}
		\begin{figure}[h]
			\centering
			\includegraphics[scale=0.55]{72.png}
			\caption{}
			\label{fig:im-72}
		\end{figure}
	\end{center}
	\subsection{Resistività}
	La resistenza di un conduttore dipende dalle sue dimensioni, dalla sua forma e dal materiale di cui è fatto. \\
	Considerando un conduttore a sezione costante:
	\begin{center}
		\begin{figure}[h]
			\centering
			\includegraphics[scale=0.55]{73.png}
			\caption{}
			\label{fig:im-73}
		\end{figure}
	\end{center}
	Si trova
	$$R = \frac{\rho h}{S} \qquad \qquad [\Omega] = \frac{V}{A} \quad (Ohm)$$
	La dipendenza della resistenza di un conduttore dal materiale è rappresentata da un fattore di proporzionalità chiamato resistività $\rho$ (si misura in $\Omega$ m).
	$$\sigma = \frac{1}{\rho} \quad \text{conducibilità}$$
	Tabella resistività dei vari metalli:
	\begin{center}
		\begin{tabular}{|c|c|}
			\hline
			\textbf{Metallo} & \textbf{$\rho$ ($\Omega \cdot \text{m}$)} \\ 
			\hline
			Argento      & $1,6 \cdot 10^{-8}$  \\
			Rame         & $1,7 \cdot 10^{-8}$  \\
			Alluminio    & $2,8 \cdot 10^{-8}$  \\
			Ottone       & $\sim 7 \cdot 10^{-8}$ \\
			Nichel       & $7,8 \cdot 10^{-8}$  \\
			Ferro        & $10 \cdot 10^{-8}$   \\
			Acciaio      & $\sim 11 \cdot 10^{-8}$ \\
			Costantana   & $49 \cdot 10^{-8}$   \\
			Nichelcromo  & $100 \cdot 10^{-8}$  \\ 
			\hline
		\end{tabular}
	\end{center}
	*Alla temperatura di 20$^\degree$\\\\
	Tabella resistività dei vari isolanti:
	\begin{center}
		\begin{tabular}{|c|c|}
			\hline
			\textbf{Isolante} & \textbf{$\rho$ ($\Omega \cdot \text{m}$)} \\ 
			\hline
			Polietilene                 & $2 \cdot 10^{11}$   \\
			Vetro                       & $\sim 10^{12}$      \\
			Porcellana non vetrificata  & $\sim 10^{12}$      \\
			Ebanite                     & $\sim 10^{13}$      \\
			Resina epossidica           & $\sim 10^{15}$      \\ 
			\hline
		\end{tabular}
	\end{center}
	La resistività nella maggior parte dei conduttori metallici puri è una funzione crescente della temperatura. \\
	In un intervallo limitato (qualche decina di gradi) intorno alla temperatura di ambiente la relazione è praticamente lineare.
	\begin{center}
		\begin{figure}[h]
			\centering
			\includegraphics[scale=0.55]{74.png}
			\caption{$\rho = \rho_0 (1 + \alpha \Delta T)$}
			\label{fig:im-74}
		\end{figure}
	\end{center}
	\begin{center}
		\begin{figure}[h]
			\centering
			\includegraphics[scale=0.55]{75.png}
			\caption{Grafico di un super-conduttore}
			\label{fig:im-75}
		\end{figure}
	\end{center}
	L’8 Aprile 1911, Kamerlingh Onnes (Leiden) misurando la resistenza elettrica del mercurio trovò che questa diventa nulla entro gli errori a temperature inferiori ai 4.2 K (la temperatura alla quale Onnes era appena riuscito a liquefare l’elio).
	\begin{center}
		\begin{figure}[h]
			\centering
			\includegraphics[scale=0.55]{76.png}
			\caption{Kamerlingh Onnes}
			\label{fig:im-76}
		\end{figure}
	\end{center}
	\subsection{Legge di Ohm in Forma Locale}
	La legge di Ohm in forma locale lega due quantità nello stesso punto.\\\\
	Se un campo elettrico $E$ viene applicato ad un materiale conduttore, nel materiale si produce una corrente di densità $J$.\\
	La densità di corrente in un punto del materiale dipende dal campo elettrico in quel punto.
	\begin{center}
		\begin{figure}[h]
			\centering
			\includegraphics[scale=0.55]{77.png}
			\caption{}
			\label{fig:im-77}
		\end{figure}
	\end{center}
	$$V_A - V_B = dV = R dI = R J dS $$
	Essendo $E$ moltiplicato per $dI$ deduciamo che passa corrente infinitesima.
	$$= \rho \frac{dl}{dS} J dS = \rho J dl$$
	$$dV = E dl = \rho J dl$$
	$$\boxed{\vec{E} = \rho \vec{J}}$$ 
	$$\boxed{\vec{J} = \sigma \vec{E}}$$
	\subsection{Modello di Drude per i Metalli}
	Il modello di Drude deve riprodurre la legge di Ohm a livello microscopico per essere valido.\\\\
	La principale caratteristica della conduzione nei metalli che un modello microscopico soddisfacente deve riprodurre è la legge di Ohm.
	$$\mathbf{j} = \sigma \mathbf{E}$$
	Poiché $\mathbf{j} = nq\mathbf{V}_d$ si ha $\mathbf{V}_d = \frac{\sigma \mathbf{E}}{nq}$.
	La legge di Ohm implica che vd è proporzionale a campo elettrico.\\\\
	L’applicazione di un campo fa muovere i portatori con una velocità media costante, ma se $qE$ fosse l’unica forza in gioco i portatori si dovrebbero muovere di moto accelerato.\\
	Gli elettroni urtano gli ioni del cristallo subendo una brusca variazione di direzione del moto.
	\begin{center}
		\begin{figure}[h]
			\centering
			\includegraphics[scale=0.55]{78.png}
			\caption{}
			\label{fig:im-78}
		\end{figure}
	\end{center}
	Quando si applica un campo elettrico $E$ gli elettroni sono sottoposti ad una forza $F = -e E$. \\
	Questa è l’unica forza che agisce tra un urto e l’altro.\\\\
	Dopo ogni collisione gli elettroni acquistano una direzione di moto casuale.\\\\
	Dalla seconda legge di Newton (assumendo $E$ diretto lungo $x$)
	$$a_x = -e \frac{E}{m}$$
	Che vediamo essere negativo poichè gli elettroni si muovono in verso opposto al campo.
	$$v_x(t) = v_x(0) - \frac{eEt}{m}$$
	Mediando sugli elettroni ($< v_x(0) >= 0$) (Perchè acquistano direzione casuale).\\\\
	Velocità media degli elettroni:
	$$< v_x > = v_d = - \frac{eE\tau}{m}$$
	$\tau$ caratterizza l’intervallo di tempo medio tra gli urti.\\\\
	Coducibilità:
	$$\sigma = \frac{ne^2\tau}{m}$$
	$v_d << v_T$ (ovvero: Velocità di deriva molto minore di velocità elettroni) $\tau$ non dipende da $E$.
	\subsection{Effetto Joule}
	Consideriamo una carica $dq$ che si muova attraversando una differenza di potenziale $\Delta V = V_A - V_B$. \\
	Il campo elettrico compie il lavoro:
	$$dL = \Delta V \, dq = \Delta V \, i \, dt$$
	Pertanto la potenza elettrica spesa è
	$$P = \frac{dL}{dt} = \Delta V \, i$$
	Se vale la legge di Ohm, si ha
	$$P = \Delta V \, i = R \, i^2 = \frac{\Delta V^2}{R}$$
	Il passaggio di corrente attraverso un conduttore metallico per un tempo $t$ comporta dunque il lavoro:
	$$L = \int R \, i^2 \, dt$$
	Questo lavoro è necessario per vincere la resistenza opposta dal reticolo cristallino al moto ordinato degli elettroni (urti).\\
	Da un punto di vista termodinamico, esso viene assorbito dal conduttore la cui energia interna aumenta.\\
	Questo lavoro compensa la perdita di energia dovuta agli urti col reticolo cristallino. \\
	Da un punto di vista termodinamico esso viene assorbito dal conduttore la cui energia interna aumenta.\\\\
	L’effetto di riscaldamento di un conduttore percorso da corrente si chiama \textbf{effetto Joule}.
	\subsection{Resistori in Serie e in Parallelo}
	Conduttori ohmici caratterizzati da un determinato valore della resistenza (alla temperatura ambiente) sono elementi molto usati nei circuiti elettrici. \\
	Essi vengono chiamati resistori.
	\begin{center}
		\begin{figure}[h]
			\centering
			\includegraphics[scale=0.55]{79.png}
			\caption{Resistore}
			\label{fig:im-79}
		\end{figure}
	\end{center}
	Più resistori possono essere collegati tra loro, tipicamente da fili o piattine metallici, la cui resistenza è trascurabile.
	\subsubsection{Resistori in Serie}
	\begin{center}
		\begin{figure}[h]
			\centering
			\includegraphics[scale=0.55]{80.png}
			\caption{Resistori in serie}
			\label{fig:im-80}
		\end{figure}
	\end{center}
	Due resistori sono collegati in serie quando hanno un estremo in comune e \textbf{la corrente che li attraversa è la stessa}. \\
	Applicando la legge di Ohm:
	$$V_A - V_B = R_1 i$$
	$$V_B - V_C = R_2 i$$
	$$V_A - V_C = (R_1 + R_2)i = R_{eq} i$$
	$$R_{eq} = R_1 + R_2$$
	\subsubsection{Resistori in Parallelo}
	\begin{center}
		\begin{figure}[h]
			\centering
			\includegraphics[scale=0.55]{81.png}
			\caption{Resistori in parallelo}
			\label{fig:im-81}
		\end{figure}
	\end{center}
	Due resistori si dicono in parallelo quando sono collegati tra loro in entrambi gli estremi e quindi la d.d.p. (Differenza di potenziale) ai loro capi è la stessa.
	$$i = i_1 + i_2 = \frac{\Delta V}{R_1} + \frac{\Delta V}{R_2} = \Delta V \left( \frac{1}{R_1} + \frac{1}{R_2} \right)$$
	$$\frac{1}{R_{eq}} = \frac{1}{R_1} + \frac{1}{R_2}$$
	Il reciproco della resistenza equivalente è uguale alla somma dei reciproci delle singole resistenze.\\
	$R_{eq}$ risulta dunque minore di ciascuna delle $R_i$.
	\subsection{Amperometri e Voltmetri}
	\begin{figure}[h]
		
		\begin{subfigure}{0.5\textwidth}
			\includegraphics[scale=0.55]{82.png} 
			\caption{Amperometro (Richiede $r_i$ piccola)}
			\label{fig:im-82}
		\end{subfigure}
		\begin{subfigure}{0.5\textwidth}
			\includegraphics[scale=0.55]{83.png}
			\caption{Voltmetro (Richiede $r_i$ grande)}
			\label{fig:im-83}
		\end{subfigure}
		
		\caption{}
		\label{fig:image7}
	\end{figure}
	\subsection{Forza Elettromotrice}
	In un conduttore gli elettroni si muovono con una velocità $v_d$ dai punti con potenziale minore ai punti con potenziale maggiore sotto l’effetto del campo elettrico presente nel conduttore.\\\\
	Il lavoro effettuato dal campo elettrico viene assorbito dagli ioni del conduttore la cui energia interna aumenta (effetto Joule).
	\begin{center}
		\begin{figure}[h]
			\centering
			\includegraphics[scale=0.55]{84.png}
			\caption{}
			\label{fig:im-84}
		\end{figure}
	\end{center}
	Una sorgente esterna di energia è necessaria per far circolare con continuità una corrente in un circuito elettrico. \\
	Un dispositivo che \textbf{mantenga una differenza di potenziale costante} tra due punti del circuito è detto \textbf{generatore} di forza elettromotrice.\\\\
	Esaminiamo il caso semplice di un conduttore collegato ai poli $A$ e $B$ di un generatore sui cui poli sono accumulate le cariche $+q$ e $-q$.
	Il campo elettrostatico $E_{el}$ prodotto da tali cariche è sempre diretto da $A$ verso $B$, sia nel conduttore che all’interno del generatore.\\
	Quindi all’interno del generatore $E_{el}$ si oppone al moto delle cariche.
	$$\oint \mathbf{E}_{el} \cdot \mathbf{ds} = 0$$
	Per ottenere nel circuito una corrente $i$ è necessaria la presenza nel generatore di un campo $E^*$ (\textbf{campo elettromotore}) che muove le cariche contro $E_{el}$ e mantiene la differenza di potenziale tra $A$ e $B$.\\\\
	La sorgente di f.e.m. (Forza ElettroMotrice) deve quindi avere al suo interno forze di natura non elettrostatica, non conservative, che possono determinare il moto continuo delle cariche.
	$$\mathbf{E}^* \qquad \qquad \text{campo elettromotore}$$
	$$\mathbf{E} = \mathbf{E}_{el} \qquad \qquad \text{nei conduttori}$$ 
	$$\mathbf{E} = \mathbf{E}_{el} + \mathbf{E}^* \qquad \qquad \text{nel generatore}$$ 
	Il dispositivo che genera il campo elettromotore, e quindi la f.e.m., può sfruttare ad es. \textbf{reazioni chimiche} (pile e accumulatori) o il \textbf{fenomeno dell’induzione elettromagnetica}.\\\\
	Pertanto la f.e.m. $f$ è il \textbf{lavoro per unità di carica} effettuato dal campo elettromotore:
	$$f = \int_{B}^{A} \mathbf{E}^* \cdot d\mathbf{s} = \oint \mathbf{E}^* \cdot d\mathbf{s}$$
	\textbf{Si misura in Volt e non è una forza}.\\\\
	Un generatore non ideale è caratterizzato anche da un resistenza interna $r$. \\
	Quindi $f$ è la d.d.p. misurata ai capi del generatore a circuito aperto.
	\begin{center}
		\begin{figure}[h]
			\centering
			\includegraphics[scale=0.55]{85.png}
			\caption{}
			\label{fig:im-85}
		\end{figure}
	\end{center}
	$$\boxed{V_A - V_B = f - I_r}$$
	A circuito aperto: $I = 0$
	$$\Downarrow$$
	$$V_A - V_B = f$$
	La f.e.m. è la differenza di potenziale tra i morsetti del generatore a circuito aperto.
	\subsection{Legge di Ohm generalizzata e $2^{a}$ legge di Kirchhoff}
	Gli elementi geometrici distintivi di un circuito sono i \textbf{nodi} e i \textbf{rami}. \\
	Un nodo è un punto nel quale convergono almeno tre conduttori.\\
	I nodi sono collegati da rami, in cui possono esserci componenti attivi (generatori) e componenti passivi (resistori).\\\\
	All’interno di una rete è possibile individuate determinati cammini chiusi, detti maglie, costituiti da più rami; un dato ramo può pertanto appartenere a più maglie.
	\begin{center}
		\begin{figure}[h]
			\centering
			\includegraphics[scale=0.55]{86.png}
			\caption{Ramo}
			\label{fig:im-86}
		\end{figure}
	\end{center}
	$$V_A - I R_1 + f_1 - I R_2 - I R_3 - f_2 - i R_4 = V_B$$
	$$V_A - V_B + \sum_{ALG} f_i = \sum_{RAMO} I R_i \qquad \qquad \text{(comprensiva delle resistenze interne)}$$
	Una conseguenza è
	\begin{center}
		\begin{figure}[h]
			\centering
			\includegraphics[scale=0.55]{87.png}
			\caption{Maglia}
			\label{fig:im-87}
		\end{figure}
	\end{center}
	Legge delle maglie (2ª di Kirchhoff):
	$$\sum_{ALG} f_i = \sum_{MAGLIA} R_i I_i$$
	\subsection{Carica di un Condensatore Attraverso un Resistore}
	\begin{center}
		\begin{figure}[h]
			\centering
			\includegraphics[scale=0.55]{88.png}
			\caption{}
			\label{fig:im-88}
		\end{figure}
	\end{center}
	Al tempo $t=0$ viene chiuso l’interruttore e il generatore inizia a prelevare cariche positive dai conduttori connessi al polo negativo e a portarle al polo positivo. \\
	Quindi sulle armature del condensatore compaiono le cariche $+q$ e $-q$. \\
	Il processo continua fino a quando la carica del condensatore raggiunge il valore massimo $q_m = Cf$, cui corrisponde la d.d.p. tra le armature $V_A - V_B = f$.
	$$f = Ri(t) + \frac{q(t)}{C}$$
	$$i(t) = \frac{dq}{dt}$$
	\begin{center}
		\begin{figure}[h]
			\centering
			\includegraphics[scale=0.55]{89.png}
			\caption{}
			\label{fig:im-89}
		\end{figure}
	\end{center}
	$$R \frac{dq}{dt} = f - \frac{q}{C}$$
	$$\frac{dq}{q - Cf} = -\frac{dt}{RC} \qquad \qquad \int_{0}^{q} \frac{dq^\prime}{q^\prime - Cf} = -\int_{0}^{t} \frac{dt^\prime}{RC}$$
	$$\ln\left(\frac{q - Cf}{-Cf}\right) = -\frac{t}{RC}$$
	$$q(t) = Cf(1 - e^{-\frac{t}{RC}})$$
	$$\tau = RC$$
	$$i(t) = \frac{dq}{dt} = \frac{f}{R} e^{-\frac{t}{RC}}$$
	$$\tau = RC$$
	\begin{center}
		\begin{figure}[h]
			\centering
			\includegraphics[scale=0.55]{90.png}
			\caption{Grafico del potenziale in funzione del tempo}
			\label{fig:im-90}
		\end{figure}
	\end{center}
	\begin{center}
		\begin{figure}[h]
			\centering
			\includegraphics[scale=0.55]{91.png}
			\caption{Grafico della corrente in funzione del tempo}
			\label{fig:im-91}
		\end{figure}
	\end{center}
	Quando si carica un condensatore connettendolo ad un generatore, la d.d.p. finale ai capi del condensatore e uguale alla f.e.m. $f$ del generatore e la carica finale è $q_m = Cf$.\\\\
	Questi valori sono raggiunti asintoticamente. \\
	La corrente nel circuito e massima nell’istante $t=0$ e decresce esponenzialmente nel tempo, con costante di tempo $\tau = RC$.
\end{document}